\chapter{현대 LLM의 구성요소와 실행 경로}
\label{chap:phase3}

\section{결론부터: 현대 모델은 하나의 고정 설계가 아니다}

``현대 LLM architecture''라는 단일 표준은 없다. 공개 모델들은 decoder-only causal
Transformer라는 공통 골격을 자주 사용하지만 normalization, position encoding,
MLP gating, attention head grouping, bias, weight tying, mixture-of-experts 여부와 context
확장법은 서로 다르다. 따라서 정확한 순전파를 재현하려면 논문만이 아니라 해당
checkpoint의 configuration과 실행 코드를 함께 확인해야 한다.

이 장은 dense Llama 계열에서 널리 알려진 다음 경로를 대표 사례로 사용한다.
\begin{equation}
  \text{RMSNorm}\to\text{GQA with RoPE}\to+
  \to\text{RMSNorm}\to\text{SwiGLU}\to+.
\end{equation}
Meta가 공개한 Llama 3 reference implementation은 attention과 feed-forward 전에
RMSNorm을 적용하고, query와 key에 rotary embedding을 적용하며, KV head를 query
head 수에 맞게 반복하고, SiLU-gated feed-forward를 계산한다
\citep{meta2024llama3code,grattafiori2024llama3}. 이는 특정 공개 구현에 관한 사실이지
모든 LLM에 관한 정의가 아니다.

\section{decoder-only와 causal attention}

decoder-only autoregressive language model은 encoder 없이 causal self-attention block을
반복하여 다음 joint distribution을 모델링한다.
\begin{equation}
  p_{\theta}(x_{1:T})=\prod_{t=1}^{T}p_{\theta}(x_t\mid x_{<t})
\end{equation}
원래 Transformer는 encoder와 decoder를 함께 사용하여 machine translation을 수행했지만
\citep{vaswani2017attention}, GPT 계열은 position $t$의 activation이 $x_{>t}$에 의존하지
않도록 미래 position을 attention에서 가린다.

``decoder-only''는 원래 encoder--decoder Transformer의 decoder를 그대로 떼었다는
뜻으로만 이해하면 부정확할 수 있다. language-model decoder에는 encoder output을
읽는 cross-attention sublayer가 없고, causal self-attention과 position-wise MLP가
주된 block을 이룬다. prompt와 생성된 token도 같은 stream에 놓인다.

causal mask는 정보 의존성을 제한하지만 계산 순서를 항상 token별 직렬로 만들지는
않는다. 훈련과 prompt prefill에서는 전체 sequence의 query를 한 번에 병렬 계산할 수
있다. 생성 단계에서는 아직 존재하지 않는 다음 token 때문에 token별 반복이 필요하다.

\section{pre-normalization과 RMSNorm}

LayerNorm은 feature 축의 평균과 분산을 사용한다.
\begin{equation}
  \operatorname{LN}(\vect{x})=
  \vect{g}\odot
  \frac{\vect{x}-\mu(\vect{x})}
       {\sqrt{\sigma^2(\vect{x})+\epsilon}}+\vect{b}.
\end{equation}
RMSNorm은 평균을 빼지 않고 root mean square로 scale을 정규화한다
\citep{zhang2019rmsnorm}.
\begin{equation}
  \operatorname{RMSNorm}(\vect{x})=
  \vect{g}\odot
  \frac{\vect{x}}
       {\sqrt{d^{-1}\sum_{i=1}^{d}x_i^2+\epsilon}}.
  \label{eq:rmsnorm}
\end{equation}
$\vect{g}$는 학습되는 coordinate-wise gain이다. 구현에 따라 bias를 두지 않는다.
\index{RMSNorm@RMSNorm}

post-norm block은 sublayer와 residual addition 뒤에 normalization을 적용한다. pre-norm
block은 sublayer가 residual stream을 읽기 전에 normalization을 적용한다. 두 배치는
함수가 다르며 checkpoint의 가중치를 그대로 교환해 사용할 수 없다. pre-norm의
초기 gradient 거동을 분석한 연구는 post-norm보다 warm-up 의존성이 낮아질 수 있는
조건을 제시했지만, 이것이 모든 학습 설정에서 항상 더 우수하다는 정리는 아니다
\citep{xiong2020layernorm}.

\begin{factbox}
RMSNorm은 residual stream 자체를 매번 unit RMS로 덮어쓰지 않는다. pre-norm 모델에서는
정규화된 사본을 attention이나 MLP가 읽고, sublayer output은 정규화되지 않은 residual
stream에 더해진다.
\end{factbox}

\section{RoPE}

Rotary Position Embedding은 짝을 이룬 query와 key 좌표를 position에 따라 회전시킨다
\citep{su2021roformer}. 한 2차원 coordinate pair에 대해
\begin{equation}
  \mat{R}_{\theta,t}=
  \begin{bmatrix}
    \cos(t\theta)&-\sin(t\theta)\\
    \sin(t\theta)& \cos(t\theta)
  \end{bmatrix}
\end{equation}
를 적용한다. 여러 frequency $\theta_r$에 대한 block-diagonal rotation이 전체
head dimension을 구성한다.
\index{RoPE@RoPE}

RoPE를 적용한 query--key 내적은 두 position의 상대적 offset을 포함한다. position
$i,j$의 회전된 query와 key를 $\tilde{\vect{q}}_i=\mat{R}_i\vect{q}_i$,
$\tilde{\vect{k}}_j=\mat{R}_j\vect{k}_j$라고 하면
\begin{equation}
  \tilde{\vect{q}}_i^{\mathsf T}\tilde{\vect{k}}_j
  =\vect{q}_i^{\mathsf T}\mat{R}_i^{\mathsf T}\mat{R}_j\vect{k}_j
  =\vect{q}_i^{\mathsf T}\mat{R}_{j-i}\vect{k}_j.
\end{equation}
따라서 attention score에 상대적 offset $j-i$가 구조적으로 들어간다. RoPE는 일반적으로
value를 회전하지 않으며, base frequency와 장문 context scaling 방식은 모델별로
확인해야 한다.

\begin{limitbox}
RoPE가 훈련 길이 밖의 모든 position에 안정적으로 일반화한다는 보장은 위 대수에서
나오지 않는다. context extension은 별도의 scaling 방법과 평가가 필요한 경험적 문제다.
\end{limitbox}

\section{SwiGLU}

GLU 계열은 두 affine projection 가운데 하나를 gate로 사용한다
\citep{shazeer2020glu}. bias를 생략한 대표적인 SwiGLU feed-forward는
\begin{equation}
  \operatorname{SwiGLU}(\vect{x})=
  \mat{W}_{\mathrm{down}}
  \left(\operatorname{SiLU}(\mat{W}_{\mathrm{gate}}\vect{x})
  \odot(\mat{W}_{\mathrm{up}}\vect{x})\right),
  \label{eq:swiglu}
\end{equation}
\begin{equation}
  \operatorname{SiLU}(a)=a\,\sigma(a)
\end{equation}
로 쓸 수 있다. 행렬 방향은 column-vector 표기다. 코드가 row-vector를 사용하면 전치된
shape로 나타난다.
\index{SwiGLU@SwiGLU}

SwiGLU에서는 gate projection의 activation이 up projection이 전달하는 내용을 조절한다.
$\mat{W}_{\mathrm{gate}},\mat{W}_{\mathrm{up}}$은 $d$에서 $d_{\mathrm{ff}}$로,
$\mat{W}_{\mathrm{down}}$은 다시 $d$로 보내며, 두 중간 벡터를 element-wise product로
결합한다.

``SwiGLU가 MLP의 지식을 저장한다''는 문장은 수학적 정의가 아니다. 수학적으로 확정된
내용은 \cref{eq:swiglu}의 함수와 parameter shape다. feature나 memory에 관한 설명은
가중치와 activation을 분석한 경험적 가설이다.

\section{MHA, MQA와 GQA}

다중 헤드 어텐션(MHA)은 query, key, value에 같은 수의 head를 둔다. query head 수를
$H_q$, key/value head 수를 $H_{kv}$라고 하면 MHA에서는 $H_q=H_{kv}$다. 다중 질의
어텐션(MQA)은 여러 query head가 하나의 key head와 value head를 공유하므로
$H_{kv}=1$이다 \citep{shazeer2019mqa}. 그룹 질의 어텐션(GQA)은 그 사이에 있으며
$1<H_{kv}<H_q$인 구성을 허용한다 \citep{ainslie2023gqa}.
\index{GQA@그룹 질의 어텐션(GQA)}

GQA는 각 query head를 하나의 shared key/value head에 대응시킨다. $H_q$가 $H_{kv}$의
배수이고 $g=H_q/H_{kv}$일 때, query head $h$가 사용하는 key/value head의 번호는
대표적으로
\begin{equation}
  r(h)=\left\lfloor\frac{h}{g}\right\rfloor
\end{equation}
로 정한다. 그러면 head $h$의 계산은
\begin{align}
  \mat{A}^{h}&=\softmax\!\left(
  \frac{\mat{Q}^{h}(\mat{K}^{r(h)})^{\mathsf T}}{\sqrt{d_h}}+\mat{M}
  \right),\\
  \mat{O}^{h}&=\mat{A}^{h}\mat{V}^{r(h)}
\end{align}
가 된다. ``key/value head를 반복한다''는 구현은 같은 $\mat{K}^{r(h)}$와
$\mat{V}^{r(h)}$를 여러 query head에 대응시키기 위한 shape 조작이다. 실제 데이터를
복사하지 않고 broadcast하거나 특화된 kernel에서 직접 공유할 수도 있다.

decode 단계에서 매 token마다 저장하는 key와 value 원소 수는 대략
$2LH_{kv}d_h$다. 따라서 $H_{kv}$를 줄이면 KV cache의 메모리 대역폭과 저장량을 줄일
수 있다. 그러나 MHA, MQA, GQA의 품질과 속도 차이는 모델 크기, 학습법, kernel과
hardware에 의존한다. GQA 논문은 MHA checkpoint를 적은 추가 학습으로 GQA로 바꾸는
방법과 그 실험 결과를 제시했지만, 모든 모델에서 같은 trade-off를 보장하는 정리는
아니다 \citep{ainslie2023gqa}.

\begin{factbox}
GQA는 query의 수를 줄이는 방법이 아니다. query head는 여러 개를 유지하되, 각 그룹이
같은 key와 value를 읽는다. 따라서 attention weight는 query head마다 다를 수 있다.
\end{factbox}

\section{tokenization과 vocabulary}

모델은 문자열을 직접 받지 않는다. tokenizer가 byte 또는 Unicode 문자열을 vocabulary의
정수 ID 열로 바꾼다. Byte Pair Encoding(BPE)은 반복적으로 빈도가 높은 인접 symbol
pair를 합치는 압축 알고리즘에서 출발했으며, 신경망 번역에서는 subword vocabulary를
만드는 방법으로 사용되었다 \citep{sennrich2016bpe}. SentencePiece는 사전 분절된 단어를
요구하지 않고 raw sentence에서 subword model을 학습하는 구현을 제안했다
\citep{kudo2018sentencepiece}.
\index{tokenizer@토크나이저(tokenizer)}

token과 단어는 일반적으로 일대일로 대응하지 않는다. tokenization 함수 $\tau$와
detokenization 함수 $\delta$를
\begin{equation}
  \tau:\mathcal{S}\to\{0,\ldots,V-1\}^{*},\qquad
  \delta:\{0,\ldots,V-1\}^{*}\to\mathcal{S}
\end{equation}
로 쓰자. 여기서 $\mathcal{S}$는 문자열 집합이고 별표는 길이가 가변적인 열을 뜻한다.
공백, 구두점, byte 조각 또는 여러
문자가 한 token이 될 수 있다. 같은 표면 문자열도 앞 공백과 normalization 방식에 따라
서로 다른 ID 열이 될 수 있다.

tokenizer는 모델 가중치와 분리된 임의의 전처리기가 아니다. embedding과 LM head의
row 의미가 vocabulary ID에 고정되어 있으므로 checkpoint와 맞는 tokenizer를 사용해야
한다. 특수 token의 ID, 문장 시작과 끝 처리, 대화 template도 실제 입력 열을 바꾼다.

\begin{limitbox}
모델이 ``몇 단어를 보았다''는 표현은 token 수를 정확히 나타내지 않는다. 계산량,
context 길이와 loss를 재현하려면 원 문자열뿐 아니라 tokenizer 버전, normalization,
special-token 설정과 최종 token IDs를 기록해야 한다.
\end{limitbox}

\section{KV cache, prefill과 decode}

길이 $T$의 prompt를 처음 처리하는 단계를 prefill이라고 한다. 모든 prompt position의
query, key와 value를 만들고 causal attention을 병렬로 계산한 뒤, 각 layer의 key와
value를 cache에 보관할 수 있다. 이어서 token 한 개를 생성하는 decode 단계에서는 새
position의 query, key와 value만 계산하고, 새 query가 과거의 cached key/value와 새
key/value를 함께 읽는다.
\index{KV cache@KV 캐시}
\index{prefill@prefill과 decode}

decode의 새 query는 과거 cache에 새 key와 value를 붙인 전체 prefix를 읽는다. layer
$\ell$의 과거 cache를 $\mat{K}^{(\ell)}_{1:T},\mat{V}^{(\ell)}_{1:T}$라고 하고, 새
token의 projection을 $\vect{q}_{T+1},\vect{k}_{T+1},\vect{v}_{T+1}$로 쓰면
\begin{align}
  \mat{K}_{1:T+1}&=\operatorname{Append}(\mat{K}_{1:T},\vect{k}_{T+1}),\\
  \mat{V}_{1:T+1}&=\operatorname{Append}(\mat{V}_{1:T},\vect{v}_{T+1}),\\
  \vect{o}_{T+1}&=
  \softmax\!\left(\frac{\vect{q}_{T+1}\mat{K}_{1:T+1}^{\mathsf T}}
  {\sqrt{d_h}}\right)\mat{V}_{1:T+1}
\end{align}
를 계산한다. 과거 token의 key와 value는 다시 계산하지 않는다.

cache가 없으면 생성 step $s$마다 앞선 $T+s$개 position의 K/V를 다시 만들어야 한다.
cache는 이 중복 projection을 없애지만, sequence 길이에 비례하는 저장 공간과 매 step의
cache 읽기 비용이 필요하다. batch의 sequence가 서로 다른 속도로 끝나면 memory block을
효율적으로 배치하는 별도 시스템 문제가 생긴다. PagedAttention은 운영체제의 가상
메모리와 비슷한 block 단위 관리로 KV cache의 낭비와 공유 문제를 다루는 시스템을
제안했다 \citep{kwon2023pagedattention}.

\begin{table}[htbp]
\centering
\caption{prefill과 한 token decode의 계산 경로}
\label{tab:prefill-decode}
\begin{tabularx}{\textwidth}{@{}lYY@{}}
\toprule
구분 & prefill & decode 한 단계 \\
\midrule
입력 position & prompt 전체 & 새 position 한 개 \\
새로 만드는 Q/K/V & prompt 전체 & 새 position만 \\
attention이 읽는 K/V & 각 position의 허용된 prefix & 과거 cache와 새 position \\
주된 병렬 축 & 많은 token position & batch, head 및 내부 행렬 연산 \\
주된 자원 병목 & 큰 행렬 연산과 메모리 & cache 대역폭과 작은 연산의 반복 \\
출력 사용 & 보통 마지막 prompt logit & 새 token의 logit \\
\bottomrule
\end{tabularx}
\end{table}

KV cache는 수학적 모델의 조건부분포를 바꾸지 않는 재사용 장치다. 다만 dtype, kernel,
양자화와 연산 순서가 바뀌면 유한정밀도 결과가 완전히 bit-wise 동일하지 않을 수 있다.

\section{FlashAttention과 수치적 동등성}

표준 attention을 순진하게 구현하면 $\mat{S}=\mat{Q}\mat{K}^{\mathsf T}$와
$\mat{A}=\softmax(\mat{S})$를 고대역폭 메모리에 저장하므로 sequence 길이에 대해
quadratic한 중간 메모리가 든다. FlashAttention은 attention 행렬을 block으로 나누고
on-chip SRAM과 고대역폭 메모리 사이의 읽기와 쓰기를 줄이는 순서를 사용한다. 원
논문은 근사 attention이 아니라 exact attention algorithm으로 제시하며, 역전파에
필요한 중간값도 재계산과 online softmax로 관리한다 \citep{dao2022flashattention}.
\index{FlashAttention@FlashAttention}

online softmax는 block별 최대값과 지수합을 재조정하여 전체 row의 softmax를 보존한다.
지금까지 처리한 score의 최대값과 지수합을 $m,l$이라고 하고 새 block의 최대값과
지수합을 $m_b,l_b$라고 할 때, 병합식은
\begin{align}
  m'&=\max(m,m_b),\\
  l'&=e^{m-m'}l+e^{m_b-m'}l_b
\end{align}
가 된다. value의 가중합도 같은 scale factor로 갱신하면 전체 score를 한꺼번에
materialize하지 않고 같은 실수식의 결과를 얻을 수 있다.

\begin{factbox}
FlashAttention은 attention의 모델 구조를 새로운 의미론으로 바꾸는 구성요소가 아니다.
같은 Q, K, V, mask와 softmax attention을 I/O 효율적인 순서로 계산하는 algorithm이다.
따라서 회로를 해석할 때는 수학적 attention과 kernel 구현을 구분해야 한다.
\end{factbox}

``exact''는 이상적인 연산식에 관한 표현이다. floating-point 덧셈은 결합법칙을 만족하지
않으므로 연산 순서가 바뀌면 마지막 bit가 달라질 수 있다. dropout, mixed precision,
kernel 버전과 hardware까지 고정하지 않으면 bit-wise 재현성을 기대해서는 안 된다.

\section{weight tying과 LM head}

weight tying은 입력 embedding과 출력 LM head가 같은 parameter를 공유하는 설계다. 두
행렬 $\mat{E},\mat{W}_U\in\R^{V\times d}$에 대해
\begin{equation}
  \mat{W}_U=\mat{E}
\end{equation}
이 등식은 두 모듈이 같은 parameter 저장소를 사용한다는 뜻이다 \citep{press2017tying}.
untied 모델은 두 행렬을
독립적으로 학습한다. tying 여부는 architecture에 따라 다르며, 같은 shape라는 사실만으로
공유를 결론낼 수 없다.

tied 모델에서도 input embedding으로서의 역할과 unembedding으로서의 역할은 계산
문맥이 다르다. 입력에서는 token ID로 row를 선택한다. 출력에서는 마지막 hidden vector와
모든 row의 내적을 한꺼번에 계산한다. final normalization이나 별도 output bias가 있으면
그 연산도 포함해야 한다.

\section{훈련이 가중치를 만드는 과정}

순전파 수식은 가중치가 어떻게 생겼는지 설명하지 않는다. autoregressive pretraining은
corpus의 token 열에서 다음 token 음의 로그가능도를 최소화한다.
\begin{equation}
  \mathcal{L}_{\mathrm{NTP}}(\theta)
  =-\sum_{t=1}^{T-1}\log p_{\theta}(x_{t+1}\mid x_{\le t}).
  \label{eq:next-token-loss}
\end{equation}
teacher forcing에서는 실제 prefix $x_{\le t}$를 모든 position에 제공하고 loss를
병렬 계산한다. 역전파가 $\nabla_{\theta}\mathcal{L}$을 구하면 optimizer가 parameter를
갱신한다. Adam은 first moment와 second raw moment의 지수이동평균을 사용하며
\citep{kingma2015adam}, AdamW는 weight decay를 gradient update에서 분리한다
\citep{loshchilov2019adamw}.

학습 결과는 architecture만으로 정해지지 않는다. 데이터의 분포와 순서, objective,
초기값, optimizer, learning-rate schedule, batch 구성, 정밀도와 분산 학습의 수치적
세부사항이 최종 parameter에 영향을 준다. 모델 크기, 데이터와 compute 사이의 경험적
scaling 관계도 보고되었지만 \citep{kaplan2020scaling,hoffmann2022chinchilla}, 이는 특정
실험 범위에 맞춘 경험 법칙이며 개별 능력이나 내부 회로를 직접 결정하는 공식이 아니다.

instruction tuning과 preference optimization은 pretraining 뒤의 출력 분포를 바꾼다.
InstructGPT는 supervised demonstrations와 사람 선호로 학습한 reward model, PPO 기반
강화학습을 조합했다 \citep{ouyang2022instructgpt,schulman2017ppo}. DPO는 명시적인
reward model과 online reinforcement-learning loop 없이 preference pair에 대한 분류형
목적함수를 최적화하는 방법을 유도했다 \citep{rafailov2023dpo}. 이들은 가능한 post-training
절차의 예이며 모든 공개 또는 상용 LLM의 동일한 recipe가 아니다.

\Needspace{9\baselineskip}
\begin{limitbox}
next-token objective를 안다고 해서 학습된 내부 알고리즘을 역으로 유일하게 복원할 수
있는 것은 아니다. 서로 다른 parameterization과 내부 표현이 같은 관측 출력 또는 매우
비슷한 loss를 만들 수 있다. 이 비식별성 때문에 내부 메커니즘에는 별도의 관찰과 개입이
필요하다.
\end{limitbox}

\section{Phase 3 학습 점검}

\begin{studybox}
\begin{enumerate}
  \item 공개 checkpoint 하나를 선택하여 layer 수, $d$, $d_{\mathrm{ff}}$, $H_q$,
        $H_{kv}$, RoPE 설정, normalization의 $\epsilon$과 weight tying 여부를 표로 만든다.
  \item $H_q=32$, $H_{kv}=8$, $d_h=128$일 때 각 query head가 공유하는 KV head를 모두
        적고, token 한 개당 layer별 cache 원소 수를 계산한다.
  \item 같은 문자열을 앞 공백 유무와 Unicode normalization을 바꾸어 token ID로 변환하고
        차이를 기록한다.
  \item 길이 $T$인 prompt의 prefill과 $n$개 token decode에서 어떤 K/V projection이
        재사용되는지 의사코드로 쓴다.
  \item FlashAttention이 줄이는 것은 수학적 연산, HBM I/O, 중간 메모리 중 무엇인지
        구분한다.
  \item pretraining objective, optimizer와 post-training objective가 각각 순전파 함수의
        어떤 parameter를 바꾸는지 설명한다.
\end{enumerate}
\end{studybox}

\begin{factbox}
Phase 3을 마치면 공개 구현의 순전파 코드를 수식과 tensor shape로 대조할 수 있어야 한다.
그러나 내부 벡터의 의미를 읽거나 특정 행동의 원인을 식별하는 일은 아직 시작하지 않았다.
\end{factbox}
